\section{Motivation: A SOTA Detector Collapses Under Event Corruptions}
\label{sec:motivation}
\advisor{point 2 -- show that published SOTA (Neftci) detection plummets under OOD input, with NO retraining; quantify the drop; hint at restoration as future work.}

The case for cheap on-chip OOD detection rests on a prior empirical claim: that these
corruptions actually break a competent detector. We test this directly and adversarially --
not by training a model to be fragile, but by taking the \emph{published, pretrained} hybrid
SNN--ANN detector~\cite{neftci2025hybrid} off the shelf and feeding it corrupted input with
\textbf{no retraining and no fine-tuning}.

\paragraph{Protocol.}
We run the released checkpoint (clean Gen1 test mAP $\approx 0.36$) through its own
inference pipeline, intercepting the event histogram \emph{after} it is loaded and
\emph{before} the backbone, and applying each Gen1-C corruption at each severity using the
same seeded generator as the benchmark. mAP (COCO AP@$[.5{:}.95]$ and AP$_{50}$) is computed
by the model's native Prophesee evaluator, so the only change versus a stock validation run
is the corruption of the input tensor. The experiment is fully reproducible from
\texttt{HybridDetection/validation\_corrupt.py} and \texttt{run\_map\_degradation.sh}; raw
numbers are summarised by \texttt{analysis/summarize\_map\_degradation.py} into
Table~\ref{tab:map-degradation}.

\begin{table}[t]
  \centering
  \todofig{Insert \texttt{results/neftci\_map\_degradation\_table.tex} once the GPU sweep is
  run. Rows = 6 corruptions, columns = severities 1--5, cells = mAP; clean baseline mAP in
  the caption. (Run on Modal/A100: needs the Gen1 test set + \texttt{gen1\_mAP36.ckpt}.)}
  \caption{Detection mAP of the pretrained hybrid detector under each corruption and
  severity (no retraining). \emph{Numbers pending the GPU sweep -- see
  Docs/paper\_figures.md once banked.}}
  \label{tab:map-degradation}
\end{table}

\paragraph{Expected reading (hypotheses to confirm with the run).}
The leaky-integrator theory (Section~\ref{sec:theory}) predicts the input-side severity of
each corruption, which transfers to detection performance: \texttt{hot\_pixel} and
\texttt{event\_flood} inject large spurious activity that should swamp object evidence and
crater mAP; \texttt{spatial\_dropout} deletes the very regions objects occupy; and
\texttt{polarity\_flip}, although nearly invisible to the membrane, inverts the ON/OFF
evidence the detection head was trained on and should degrade mAP even though our membrane
detector cannot see it -- a deliberate, instructive mismatch between \emph{detectability} and
\emph{harm}. Quantifying these drops converts ``OOD detection is nice to have'' into ``the
deployed detector fails on inputs it cannot flag.''

\paragraph{Toward restoration (future work).}
\advisor{not required for this paper, but hinting at restoration in the discussion is important.}
If a frame can be flagged as corrupted cheaply on-chip, the natural next step is to act on
the flag -- e.g.\ gating the prediction, switching to a corruption-robust head, or applying a
lightweight input correction (polarity re-balancing for \texttt{polarity\_flip}, adaptive
gain for \texttt{event\_flood}; cf.\ the biological mechanisms in
Section~\ref{sec:discussion}). We sketch this as future work rather than evaluate it here, to
keep the present contribution -- detection -- cleanly scoped.
